%% LyX 2.4.4 created this file.  For more info, see https://www.lyx.org/.
%% Do not edit unless you really know what you are doing.
\documentclass[english]{article}
\usepackage[T1]{fontenc}
\usepackage[latin9]{inputenc}
\usepackage{units}
\usepackage{enumitem}
\usepackage{amsmath}
\usepackage{amsthm}
\usepackage{amssymb}
\usepackage{geometry}
\geometry{verbose,tmargin=1in,bmargin=1in,lmargin=1in,rmargin=1in}

\makeatletter
%%%%%%%%%%%%%%%%%%%%%%%%%%%%%% Textclass specific LaTeX commands.
\numberwithin{equation}{section}
\numberwithin{figure}{section}
\newlength{\lyxlabelwidth}      % auxiliary length 

%%%%%%%%%%%%%%%%%%%%%%%%%%%%%% User specified LaTeX commands.
\usepackage{fancyhdr}
\usepackage{lastpage}
\pagestyle{fancy}
\usepackage{enumitem}
\usepackage{ifthen}
\everymath=\expandafter{\the\everymath\displaystyle}
%\DeclareMathSizes{10}{10}{10}{10}
\usepackage{multicol}

\usepackage{tikz}
\usepackage{tikz-3dplot}
\usetikzlibrary{calc,arrows}

\usetikzlibrary{patterns}
\usetikzlibrary{plotmarks}
\usepackage{pgfplots}
\pgfplotsset{compat=newest}

\usepackage{calc}
\usepackage[nomessages]{fp}

\usepackage{datetime}
% Added by lyx2lyx
\setlength{\parskip}{\smallskipamount}
\setlength{\parindent}{0pt}

\makeatother

\usepackage{babel}
\begin{document}
\renewcommand{\author}{Dr.\,James Ashe}

\newcommand{\classNum}{336}

\newcommand{\classSec}{A}

\newcommand{\examType}{HW 1.2 Solutions}

\renewcommand{\date}{\formatdate{6}{9}{2013}}

%%First page's header%%%%%%%%%%%%%%%%%%%%%%
\fancypagestyle{firststyle} %%header settings for first pagr
{
	\fancyhead[L]{MTH \classNum{}(\classSec{}) \author{}\\
	\textbf{\examType{}} ~\date{}}
	\fancyhead[C]{}
	\fancyhead[R]{\textbf{Name:\hspace{4cm}}}
	\setlength{\headsep}{14pt}
}
%%%%%%%%%%%%%%%%%%%%%%%%%%%%%%%%%%%%%%%%%%

%%Next page's header%%%%%%%%%%%%%%%%%%%%%%%
\setlist{nolistsep}
\lhead{\classNum{}(\classSec{})} 
\chead{\textbf{Name:}} 
\cfoot{\thepage{}~of~\pageref{LastPage}} 
\setlength{\headsep}{5pt} 
%%%%%%%%%%%%%%%%%%%%%%%%%%%%%%%%%%%%%%%%%%%

%%Various settings; see the preamble for other settings%%%%%
%\setenumerate[0]{leftmargin=12pt,itemindent=0pt,labelwidth=0pt}
\setlist{topsep=.5pt}
\renewcommand{\labelenumi}{\textbf{\arabic{enumi}.}}
\renewcommand{\labelenumii}{\textbf{\alph{enumii}.}}
\thispagestyle{firststyle}
%%%%%%%%%%%%%%%%%%%%%%%%%%%%%%%%%%%%%%%%%%

\newcommand{\xmax}{0}
\newcommand{\xmin}{-\xmax}
\newcommand{\xstep}{0}
\newcommand{\ymax}{0}
\newcommand{\ymin}{-\ymax}
\newcommand{\ystep}{0} 
\newcommand{\dspace}{\vspace{1cm}}
\newcommand{\xa}{0}
\newcommand{\xb}{0}
\newcommand{\ya}{1}
\newcommand{\yb}{1}
\begin{enumerate}[resume]
\item $\mathbf{x\cdot}\mathbf{y}=(x_{1},x_{2},\dots,x_{n})\cdot(y_{1},y_{2},\dots,y_{n})=x_{1}y_{1}+x_{2}y_{2}+\cdots+x_{n}y_{n}$
and $\mathbf{x\cdot}\mathbf{y}=\Vert\mathbf{x}\Vert\Vert\mathbf{y}\Vert\cos\theta$.

\begin{enumerate}[resume]
\item [\textbf{b.}]\setcounter{enumii}{2}$(2,1)\cdot(-1,1)=2(-1)+1(1)=-2+1=-1$.
$\Vert(2,1)\Vert=\sqrt{2^{2}+1^{2}}=\sqrt{5}$ and $\Vert(-1,1)\Vert=\sqrt{(-1)^{2}+1^{2}}=\sqrt{2}$.
So then $\theta=\cos^{-1}\left(\nicefrac{-1}{\sqrt{10}}\right)$ ($\approx108.435^{\circ}$). 
\item $(1,8)\cdot(7,4)=-25$. $\theta=\cos^{-1}\left(\nicefrac{-25}{\Vert(1,8)\Vert\Vert(7,4)\Vert}\right)=\cos^{-1}\left(\nicefrac{-25}{\sqrt{65}\sqrt{65}}\right)=\cos^{-1}\left(-\nicefrac{5}{13}\right)$
($\approx112.62^{\circ}$).
\item $(1,4,-3)\cdot(5,1,3)=1(5)+4(1)+(-3)(3)=5+4-9=0$. So the vectors
are orthogonal, i.e., $\theta=\nicefrac{\pi}{2}=90^{\circ}$.
\end{enumerate}
\item $\mbox{proj}_{\mathbf{y}}\mathbf{x}=\frac{\mathbf{x\cdot y}}{\Vert\mathbf{y}\Vert^{2}}\mathbf{y}$
and $\mbox{proj}_{\mathbf{x}}\mathbf{y}=\frac{\mathbf{x\cdot y}}{\Vert\mathbf{x}\Vert^{2}}\mathbf{x}$.

\begin{enumerate}[resume]
\item [\textbf{b.}]Using results from above,\\
$\mbox{proj}_{\mathbf{y}}\mathbf{x}=\frac{-1}{(\sqrt{2})^{2}}(-1,1)=-\frac{1}{2}(-1,1)$
and $\mbox{proj}_{\mathbf{x}}\mathbf{y}=\frac{-1}{(\sqrt{5})^{2}}(2,1)=-\frac{1}{5}(2,1)$.
\item [\textbf{d.}]The vectors are orthogonal, so $\mbox{proj}_{\mathbf{y}}\mathbf{x}=\mbox{proj}_{\mathbf{x}}\mathbf{y}=0$. 
\end{enumerate}
\item [\textbf{4.}]Note that scaling and position does not affect the problem.
So without loss of generality we can assume that vectors $(0,0,0),(1,0,0),(0,1,0),(1,1,0),(0,0,1),(1,0,1),(0,1,1)$,
and $(1,1,1)$ represent the eight vertices of the cube. A long diagonal
and a face angle are given by $(1,1,1)$ and $(1,0,0)$ respectively.
$(1,1,1)\cdot(1,1,0)=2$, $\Vert(1,1,1)\Vert=\sqrt{3}$, and $\Vert(1,1,0)\Vert=\sqrt{2}$.
This gives $\cos^{-1}\left(\nicefrac{2}{\sqrt{6}}\right)\approx35.3^{\circ}$.\\
\tdplotsetmaincoords{60}{115}% %%specifes orientation
%%\tdplotsetrotatedcoords{0}{20}{0}% %%specifes angle of rotation
\begin{tikzpicture} 		
	[tdplot_main_coords,
		scale=2,
		cube/.style={purple}, 
		grid/.style={very thin,gray}, 
		axis/.style={->,blue,thick},
		vector/.style={very thick,red,-stealth},
		vector guide/.style={dashed,red,thick},
		angle/.style={black,thick}]
	
	%draw a grid in the x-y plane 	
	\foreach \x in {-0.25,0,...,1.25} 		
	\foreach \y in {-0.25,0,...,1.25} 		
	{ 			
		%\draw[grid] (\x,-0.25) -- (\x,1.25); 			
		%\draw[grid] (-0.25,\y) -- (1.25,\y); 		
	} 		
	
	
	
	%draw the axes 	
	\draw[axis] (0,0,0) -- (1.5,0,0) node[anchor=west]{$x$}; 	
	\draw[axis] (0,0,0) -- (0,1.5,0) node[anchor=west]{$y$}; 	
	\draw[axis] (0,0,0) -- (0,0,1.5) node[anchor=west]{$z$};
	
	%draw the top and bottom of the cube 	
	\draw[cube] (0,0,0) -- (0,1,0) -- (1,1,0) -- (1,0,0) -- cycle; 	
	\draw[cube] (0,0,1) -- (0,1,1) -- (1,1,1) -- (1,0,1) -- cycle; 
	
	%draw the edges of the cube 	
	\draw[cube] (0,0,0) -- (0,0,1); 	
	\draw[cube] (0,1,0) -- (0,1,1); 	
	\draw[cube] (1,0,0) -- (1,0,1); 	
	\draw[cube] (1,1,0) -- (1,1,1);

	%standard tikz coordinate definition using x, y, z coords 	
	\coordinate (O) at (0,0,0);
	%tikz-3dplot coordinate definition using r, theta, phi coords 	
	\tdplotsetcoord{P}{.8}{55}{60}

	%draw vector
	%\draw[vector] (O) -- (P);
	%\draw[vector guide] (O) -- (Pxy); \draw[vector guide] (Pxy) -- (P);
	\draw[->, very thick,red,-stealth] (0,0,0) -- (1,1,1);
	\draw[->,very thick,blue,-stealth] (0,0,0) -- (1,1,0);
	%draw an arc illustrating the angle defining the orientation 
	\tdplotsetthetaplanecoords{45}	
	\tdplotdrawarc[tdplot_rotated_coords,angle]{(O)}{.35}{54.7}{90}{anchor=west}{$\phi$} %{center}{radius}{start}{end angle}
\end{tikzpicture}
	\tdplotsetmaincoords{90}{0}% %%specifes orientation
%%\tdplotsetrotatedcoords{0}{20}{0}% %%specifes angle of rotation
\begin{tikzpicture} 		
	[tdplot_main_coords,
		scale=2,
		cube/.style={purple}, 
		grid/.style={very thin,gray}, 
		axis/.style={->,blue,thick},
		vector/.style={very thick,red,-stealth},
		vector guide/.style={dashed,red,thick},
		angle/.style={black,thick}]
	
	%draw a grid in the x-y plane 	
	\foreach \x in {-0.25,0,...,1.25} 		
	\foreach \y in {-0.25,0,...,1.25} 		
	{ 			
		%\draw[grid] (\x,-0.25) -- (\x,1.25); 			
		%\draw[grid] (-0.25,\y) -- (1.25,\y); 		
	} 		
	
	
	
	%draw the axes 	
	\draw[axis] (0,0,0) -- (1.5,0,0) node[anchor=west]{$x$}; 	
	\draw[axis] (0,0,0) -- (0,1.5,0) node[anchor=west]{$y$}; 	
	\draw[axis] (0,0,0) -- (0,0,1.5) node[anchor=west]{$z$};
	
	%draw the top and bottom of the cube 	
	\draw[cube] (0,0,0) -- (0,1,0) -- (1,1,0) -- (1,0,0) -- cycle; 	
	\draw[cube] (0,0,1) -- (0,1,1) -- (1,1,1) -- (1,0,1) -- cycle; 
	
	%draw the edges of the cube 	
	\draw[cube] (0,0,0) -- (0,0,1); 	
	\draw[cube] (0,1,0) -- (0,1,1); 	
	\draw[cube] (1,0,0) -- (1,0,1); 	
	\draw[cube] (1,1,0) -- (1,1,1);

	%standard tikz coordinate definition using x, y, z coords 	
	\coordinate (O) at (0,0,0);
	%tikz-3dplot coordinate definition using r, theta, phi coords 	
	\tdplotsetcoord{P}{.8}{55}{60}

	%draw vector
	%\draw[vector] (O) -- (P);
	%\draw[vector guide] (O) -- (Pxy); \draw[vector guide] (Pxy) -- (P);
	\draw[->, very thick,red,-stealth] (0,0,0) -- (1,1,1);
	\draw[->,very thick,blue,-stealth] (0,0,0) -- (1,1,0);
	%draw an arc illustrating the angle defining the orientation 
	\tdplotsetthetaplanecoords{45}	
	\tdplotdrawarc[tdplot_rotated_coords,angle]{(O)}{.35}{54.7}{90}{anchor=west}{$\phi$} %{center}{radius}{start}{end angle}

\end{tikzpicture}
	
\item [\textbf{7.}]The vectors $2\mathbf{x}+3\mathbf{y}$ and $\mathbf{x}-\mathbf{y}$
are orthogonal if and only if their dot product is $0$. The dot product
is distributive and $\mathbf{x\cdot x=\Vert}\mathbf{x}\Vert^{2}$
so then 
\begin{eqnarray*}
(2\mathbf{x}+3\mathbf{y})\cdot(\mathbf{x}-\mathbf{y}) & = & 2\mathbf{x\cdot x}-2\mathbf{x\cdot}\mathbf{y}+3\mathbf{y}\cdot\mathbf{x}-3\mathbf{y}\cdot\mathbf{y}\\
 & = & 2\Vert\mathbf{x}\Vert^{2}+\mathbf{y\cdot}\mathbf{x}-3\Vert\mathbf{y}\Vert^{2}\\
 & = & 2(2)+\mathbf{y\cdot x}-3\mbox{.}
\end{eqnarray*}
The angle between $\mathbf{x}$ and $\mathbf{y}$ is $\nicefrac{3\pi}{4}=\cos^{-1}\left(\nicefrac{\mathbf{x\cdot}\mathbf{y}}{\sqrt{2}}\right)\Rightarrow\mathbf{x\cdot y}=-1$
(solve for $\mathbf{x\cdot y}$). Thus 
\[
2(2)+\mathbf{y\cdot x}-3=4-1-3=0
\]
 Which tells us that $2\mathbf{x}+3\mathbf{y}$ and $\mathbf{x}-\mathbf{y}$
are orthogonal.
\item [\textbf{11.}]$\mathbf{x}$ is orthogonal to each $\mathbf{v_{1}},\mathbf{v_{2}},\dots,\mathbf{v_{k}}$
implies that $\mathbf{x}\cdot\mathbf{v_{i}}=0$ for each $i=1,2,\dots,k$.
So then by the distributive property of the dot product, 
\begin{eqnarray*}
\mathbf{x}\cdot(c_{1}\mathbf{v_{1}}+c_{2}\mathbf{v_{2}}+\cdots+c_{k}\mathbf{v_{k}}) & = & c_{1}\mathbf{x\cdot}\mathbf{v_{1}}+c_{2}\mathbf{x\cdot}\mathbf{v_{2}}+\cdots+c_{k}\mathbf{x\cdot}\mathbf{v_{k}}\\
 & = & c_{1}(0)+c_{2}(0)+\cdots+c_{k}(0)\\
 & = & 0+0+\cdots+0=0\mbox{.}
\end{eqnarray*}
So then $\mathbf{x}$ is orthogonal to $c_{1}\mathbf{v_{1}}+c_{2}\mathbf{v_{2}}+\cdots+c_{k}\mathbf{v_{k}}$
for any $c_{1},c_{2},\dots,c_{k}\in\mathbb{R}$.
\item [\textbf{14}] Let $A$, $B$ be vectors as illustrated below.\\
\tdplotsetmaincoords{0}{0}% %%specifes orientation
%%\tdplotsetrotatedcoords{0}{20}{0}% %%specifes angle of rotation
\begin{tikzpicture} 		
	[tdplot_main_coords,
		scale=1.5,
		grid/.style={very thin,gray}, 
		axis/.style={->,black,thick},
		vector/.style={very thick,red,-stealth},
		vector guide/.style={dashed,red,thick},
		angle/.style={black,thick}]
	
	%draw a grid in the x-y plane 	
	%\foreach \x in {-1,0,...,1} 		
	%\foreach \y in {-1,0,...,1} 		
	{ 			
		%\draw[grid] (\x,-0.25) -- (\x,1.25); 			
		%\draw[grid] (-0.25,\y) -- (1.25,\y); 		
	} 		
	
	
	
	%draw the axes 	
	\draw[axis] (-.25,0) -- (2,0) node[anchor=west]{$x$}; 	
	\draw[axis] (0,-.25) -- (0,2) node[anchor=west]{$y$}; 	

	%standard tikz coordinate definition using x, y, z coords 	
	\coordinate (O) at (0,0);

	%draw vector
	\draw[->, very thick,red,-stealth] (0,0) -- (1.75,.5) node[anchor=south west]{$A$} node[anchor=north east,pos=.6]{$b$};
	\draw[->,very thick,red,-stealth] (0,0) -- (.4,1.5) node[anchor=south west]{$B$} node[anchor=north east,pos=.6]{$a$};
	\draw[very thick,blue,stealth-] (1.75,.5) -- (.4,1.5) node[anchor=south west,pos=.65]{$A-B$} node[anchor=north east,pos=.6]{$c$};
	\draw[black] (0,0) -- (0,0) node[anchor=south west]{$\theta$};
	%draw an arc illustrating the angle defining the orientation 
	%\tdplotsetthetaplanecoords{45}	
	%\tdplotdrawarc[tdplot_rotated_coords,angle]{(O)}{.35}{54.7}{90}{anchor=west}{$\phi$} %{center}{radius}{start}{end angle}
\end{tikzpicture}
	\\
 $b=\Vert A\Vert$, $a=\Vert B\Vert$, $c=\Vert A-B\Vert$, and $\Vert A\Vert\Vert B\Vert\cos\theta=A\cdot B$.
So then
\begin{eqnarray*}
a^{2}+b^{2}-2ab(\cos\theta^{2}) & = & \Vert B\Vert^{2}+\Vert A\Vert^{2}-2\Vert B\Vert\Vert A\Vert\left(\frac{A\cdot B}{\Vert A\Vert\Vert B\Vert}\right)\\
 & = & \Vert B\Vert^{2}+\Vert A\Vert^{2}-2A\cdot B\\
 & = & B\cdot B+A\cdot A-2A\cdot B\\
 & = & (B-A)^{2}=(A-B)^{2}=\Vert A-B\Vert^{2}\\
 & = & c^{2}
\end{eqnarray*}
($c$ is the length between $A$ and $B$ so then $c=\Vert A-B\Vert$). 
\item [\textbf{18.}] The Cauchy-Schwarz inequality says that $|\mathbf{x\cdot y}|\leq\Vert\mathbf{x}\Vert\Vert\mathbf{y}\Vert$.
Following the hint we use 
\begin{eqnarray*}
\Vert\mathbf{x}+\mathbf{y}\Vert^{2} & = & (\mathbf{x}+\mathbf{y})\cdot(\mathbf{x}+\mathbf{y})=\Vert\mathbf{x}\Vert^{2}+2\mathbf{x}\cdot\mathbf{y}+\Vert\mathbf{y}\Vert^{2}\\
 & \leq & \Vert\mathbf{x}\Vert^{2}+2|\mathbf{x\cdot y}|+\Vert\mathbf{y}\Vert^{2}\\
 & \leq & \Vert\mathbf{x}\Vert^{2}+2\Vert\mathbf{x}\Vert\Vert\mathbf{y}\Vert+\Vert\mathbf{y}\Vert^{2}\mbox{ (by Cauchy-Schwarz)}\\
 & \leq & (\Vert\mathbf{x}\Vert+\Vert\mathbf{y}\Vert)^{2}\mbox{ (factor)}
\end{eqnarray*}
\end{enumerate}

\end{document}
